% Generated by scripts/assemble.py; corrections belong in source/corrections.json.
\examhead{90}{苏州大学 2026 年数学分析试卷}{88}
% p089-b22; source PDF page 89
\blockanchor{p089-b22}
\question{1}(20 分) 解答如下问题:\par
% p089-b23; source PDF page 89
\blockanchor{p089-b23}
(1) 计算极限 $\lim_{n\to \infty}\lim_{x\to 0}\frac{\sin(nx) - n\sin x}{n^3x^3}.$\par
% p089-b24; source PDF page 89
\blockanchor{p089-b24}
(2) 设函数 $f(x) = \begin{cases} (1 + x)^{\frac{1}{x}}, & x \neq 0, \\ e, & x = 0. \end{cases}$ 求 $f'(0)$.\par
% p089-b25; source PDF page 89
\blockanchor{p089-b25}
\question{2}(20 分) 已知函数 $f$ 在 $[0,1]$ 上连续, 在 $(0,1)$ 上可微, 且 $f(0) = f(1) = 0$, $f\left(\frac{1}{2}\right) = 1$.\par
% p090-b00; source PDF page 90
\blockanchor{p090-b00}
(1) 证明: 存在 $\eta \in \left(\frac{1}{2}, 1\right)$, 使得 $f(\eta) = \eta$.\par
% p090-b01; source PDF page 90
\blockanchor{p090-b01}
(2) 证明: 存在 $\xi \in (0,1)$, 使得 $f'(\xi) - \lambda[f(\xi) - \xi] = 1$.\par
% p090-b02; source PDF page 90
\blockanchor{p090-b02}
\question{3}(15 分) 证明: 在 $\left[0, \frac{\pi}{2}\right]$ 上, 有 $\sin x \geq \frac{2}{\pi} x + \frac{1}{12\pi} x (\pi^2 - 4x^2)$.\par
% p090-b03; source PDF page 90
\blockanchor{p090-b03}
\question{4}(15 分) 解答如下问题:\par
% p090-b04; source PDF page 90
\blockanchor{p090-b04}
(1) 设 $f(x)$ 连续, 证明: $\lim_{n\to \infty}\int_0^\pi f(x)|\sin nx|\mathrm{d}x = \frac{2}{\pi}\int_0^\pi f(x)\mathrm{d}x.$\par
% p090-b05; source PDF page 90
\blockanchor{p090-b05}
(2) 设 $f(x)$ 连续, 证明: $\lim_{n\to \infty}\left(\int_{a}^{b}|f(x)|^{n}\mathrm{d}x\right)^{\frac{1}{n}} = \max_{x\in [a,b]}|f(x)|.$\par
% p090-b06; source PDF page 90
\blockanchor{p090-b06}
\question{5}(15 分) 解答如下问题:\par
% p090-b07; source PDF page 90
\blockanchor{p090-b07}
(1) 判断函数列 $f_{n}(x) = nxe^{-nx^{2}}$ 在 $(- \infty, +\infty)$ 上的一致收敛性.\par
% p090-b08; source PDF page 90
\blockanchor{p090-b08}
(2) 判断函数项级数 $\sum_{n=1}^{\infty} \frac{\ln(1 + n^2 x^2)}{n^2}$ 在 $(-2, 2)$ 上的一致收敛性.\par
% p090-b09; source PDF page 90
\blockanchor{p090-b09}
\question{6}(20 分) 解答如下问题:\par
% p090-b10; source PDF page 90
\blockanchor{p090-b10}
(1) 求积分  $\int_{0}^{1}(\ln x)^{2026}\mathrm{d}x.$\par
% p090-b11; source PDF page 90
\blockanchor{p090-b11}
(2) 求积分  $\int_{0}^{\frac{\pi}{2}}\ln(\sin x)\mathrm{d}x.$\par
% p090-b12; source PDF page 90
\blockanchor{p090-b12}
\question{7}(15 分) 解答如下问题:\par
% p090-b13; source PDF page 90
\blockanchor{p090-b13}
(1) 求 $\int_{L} (x^{2} - y) \, \mathrm{d}x + (y^{2} - x) \, \mathrm{d}y$, 其中 $L$ 是 $x^{2} + y^{2} = 4$ 的上半部分, 方向从 $(-2,0)$ 到 $(2,0)$.\par
% p090-b14; source PDF page 90
\blockanchor{p090-b14}
(2) 求由曲面 $z = \sqrt{2 - x^2 - y^2}$ 与抛物面 $z = x^2 + y^2$ 所围立体部分的表面积.\par
% p090-b15; source PDF page 90
\blockanchor{p090-b15}
\question{8}(15 分) 设 $f(x, y, z) = 2(x^3 + y^3 + z^3) - 5(x^2 + y^2 + z^2)$, 其中 $x, y, z \geq 0$, 求 $f$ 在条件 $x + y + z = 3$ 下的最值.\par
% p090-b16; source PDF page 90
\par\Needspace{6\baselineskip}
\blockanchor{p090-b16}
\question{9}(15 分) 二元函数 $u(x, y)$ 在区域 $D = [0, 1] \times [0, 1]$ 上具有二阶连续偏导数, 且在边界 $\partial D$ 上 $u = 0$. 如果\par
\nopagebreak[4]
% p090-b17; source PDF page 90
\blockanchor{p090-b17}
\displayformula{\frac {\partial^ {2} u}{\partial x ^ {2}} (x, y) + \frac {\partial^ {2} u}{\partial y ^ {2}} (x, y) = - \lambda u (x, y), \forall (x, y) \in D.}
% p090-b18; source PDF page 90
\blockanchor{p090-b18}
(1) 证明: \displayformula{\iint_{D} \left[ \left( \frac{\partial u}{\partial x} \right)^{2} + \left( \frac{\partial u}{\partial y} \right)^{2} \right] \mathrm{d}x \mathrm{d}y = \lambda \iint_{D} u^{2} \mathrm{d}x \mathrm{d}y.}\par
% p090-b19; source PDF page 90
\blockanchor{p090-b19}
(2) 若存在函数 $f, g$, 使得 $u$ 可以表示为 $u(x, y) = f(x)g(y)$, 且在区域 $D$ 内部 $u > 0$, 求 $\lambda$ 的值.\par
